\section{One Mellin integration: an elementary block budget (NS-61)}
\label{sec:ns61-smoothed-nb}

\paragraph{Status and prior art.}
The $q=2$ weighted Nyman--Beurling problem already appears in
\cite{Ehm2024}; it is not a new RH criterion. We prove its equivalence
here with the tail space stated, and use adjacent differences to obtain
an unconditional $N^{-2}$ block budget. The numerator is transformed
along with the norm. No asymptotic lower correlation estimate, NB
convergence, G2 or RH proof follows from this upper budget.

Let $H=L^2(0,\infty;dx)$ and retain $\rho_n$, $\chi$, $V_N$ and
$\kappa=\log(2\pi)-\gamma$ from Section~\ref{sec:ns53-nb-blocks}.
Define the bounded Mellin convolution
\[
 (Jf)(x)=\int_0^1f(x/u)\,\frac{du}{u},\qquad
 \sigma_n=J\rho_n,\qquad \tau=J\chi=(-\log x)\mathbf1_{(0,1)}.
\]
On the Mellin line $\Re s=1/2$, $J$ has multiplier $1/s$, so
$\norm J=2$ and $J$ is injective. Set
\[
 W_N=\operatorname{span}\{\sigma_1,\ldots,\sigma_N\},\quad
 R_N=\tau-\Pi_N\tau,\quad E_N^{(2)}=\norm{R_N}^2,
\]
where $\Pi_N$ is the orthogonal projection onto $W_N$.
The superscript is a smoothing label; $E_N^{(2)}$ is already a
squared norm, and $\norm\tau^2=2$.

\begin{proposition}[The smoothed target still detects an off-line zero]
\label{prop:ns61-rh-equivalence}
Every finite $E_N^{(2)}$ is positive, and
\[
 E_N^{(2)}\longrightarrow0\quad\Longleftrightarrow\quad\mathrm{RH}.
\]
If $\zeta(\beta+i\gamma_*)=0$ with $1/2<\beta<1$, put
$z=\beta+i\gamma_*$. Then, for every $N$,
\begin{equation}
 E_N^{(2)}\geq
 \frac{|z|^{-4}}{(2\beta-1)^{-1}+|1-z|^{-2}}>0.
 \label{eq:ns61-zero-obstruction}
\end{equation}
\end{proposition}
\begin{proof}
The original strong NB criterion and boundedness of $J$ give
$E_N^{(2)}\leq4E_N$, proving the forward approximation under RH.
For the converse, the functions $\tau$ and $\sigma_n$ belong to the
closed space
\[
 \mathcal W=L^2(0,1)\oplus
       \operatorname{span}\{x^{-1}\mathbf1_{(1,\infty)}\}.
\]
Indeed $\sigma_n(x)=1/(nx)$ for $x>1$. On this space the functional
\[
 L_z(f)=\int_0^1 f(x)x^{z-1}\,dx+\frac{c}{1-z},
 \qquad f(x)=c/x\quad(x>1),
\]
is bounded with norm squared $(2\beta-1)^{-1}+|1-z|^{-2}$.
The absolutely convergent Mellin identities on $0<\Re s<1$ are
\[
 \mathcal M\sigma_n(s)=-\frac{n^{-s}\zeta(s)}{s^2},
 \qquad \mathcal M\tau(s)=\frac1{s^2}.
\]
Thus $L_z(\sigma_n)=0$ and $L_z(\tau)=z^{-2}$, which proves
\eqref{eq:ns61-zero-obstruction}. If RH fails, the functional equation
supplies a zero on the right of the critical line. This proves the
converse without assuming a bounded inverse for $J$.
Injectivity of $J$ also transfers finite linear independence and
$\chi\notin V_N$ from the original finite problem, so $E_N^{(2)}>0$.
\end{proof}

\begin{proposition}[A unitary average identity and an elementary trace budget]
\label{prop:ns61-average-budget}
Put $U=J-I$ and $h=U\rho_1$. The operator $U$ is unitary on $H$,
$\norm h^2=\kappa$, and $h$ is nonnegative and supported in $(0,1]$.
For $n\geq2$ define
\[
 \epsilon_n=\sigma_n-\frac{n-1}{n}\sigma_{n-1}=J\delta_n.
\]
Then, as a Bochner integral in $H$,
\begin{equation}
 n\epsilon_n=\int_{n-1}^{n}h(u\,\cdot)\,du
       =U\int_{n-1}^{n}\rho_1(u\,\cdot)\,du,
 \qquad
 \norm{\epsilon_n}^2\leq
       \frac{\kappa}{n^2}\log\frac n{n-1}.
 \label{eq:ns61-average}
\end{equation}
In particular $\epsilon_n$ is nonnegative and supported in
$(0,1/(n-1)]$, with
\begin{equation}
 \int_0^\infty x\epsilon_n(x)\,dx
             =\frac{\pi^2}{24n^2(n-1)}.
 \label{eq:ns61-positive-first-moment}
\end{equation}
If $\mathcal D_N^{(2)}$ is the raw Gram of
$\epsilon_{N+1},\ldots,\epsilon_{2N}$, then
\begin{equation}
 \norm{\mathcal D_N^{(2)}}_{\rm op}
 \leq\operatorname{tr}\mathcal D_N^{(2)}\leq \mathfrak b_N,
 \qquad \mathfrak b_N=\kappa\sum_{n=N+1}^{2N}
          \frac{\log(n/(n-1))}{n^2}
 \leq\frac{3\kappa}{8N^2}.
 \label{eq:ns61-trace-budget}
\end{equation}
Moreover $N^2\operatorname{tr}\mathcal D_N^{(2)}\to3\kappa/8$.
These are unconditional analytic bounds and a trace asymptotic, not
an asymptotic formula for the residual gain.
\end{proposition}
\begin{proof}
The Mellin multiplier of $U$ is $(1-s)/s$, which has modulus one on
the critical line and is nonzero there. Hence $U$ is unitary.
Writing
\[
 A(z)=\int_0^z\frac{\{v\}}v\,dv,
 \qquad \sigma_u(x)=A(1/(ux))\quad(u>0),
\]
gives $u\partial_u\sigma_u=-\rho_u$ almost everywhere. Therefore
$\partial_u(u\sigma_u)=\sigma_u-\rho_u=h(u\,\cdot)$.
On compact positive $u$ intervals this is also the strong integral
identity: dilations are strongly continuous in $H$, and integrating
the displayed derivative recovers the pointwise formula and an
$H$-integrable majorant. Since $\sigma_1(x)=\rho_1(x)=1/x$ for
$x>1$, $h$ has the asserted support.
For $z\in[k,k+1)$ with $k\geq1$, integration in the definition of
$A$ gives $A(z)=z-k\log z+\log(k!)$, hence
$h(1/z)=k-k\log z+\log(k!)>0$.
Indeed monotonicity of $\log u$ gives
$\int_1^{k+1}\log u\,du\leq\log((k+1)!)$, or
$\log(k!)\geq k\log(k+1)-k$; use $z<k+1$.
Positivity of $\epsilon_n$ follows from its average representation.
The compactly supported $h$ has Mellin transform
$-(1-s)\zeta(s)/s^2$, which extends analytically to $\Re s>0$.
At $s=2$ it gives $\int xh(x)\,dx=\pi^2/24$. Integrating the
average identity against $x\,dx$ proves
\eqref{eq:ns61-positive-first-moment}.
The Cauchy--Schwarz inequality for a unit-length Bochner integral gives
\[
 n^2\norm{\epsilon_n}^2\leq
 \int_{n-1}^{n}\norm{h(u\,\cdot)}^2\,du
 =\kappa\int_{n-1}^{n}\frac{du}{u}.
\]
For $n-1\leq u\leq n$, $1/(n^2u)\leq u^{-3}$, whence
\[
 \mathfrak b_N\leq\kappa\int_N^{2N}u^{-3}\,du
       =\frac{3\kappa}{8N^2}.
\]
Finally, dilation continuity gives
$n^3\norm{\epsilon_n}^2\to\kappa$: rescale $x=y/n$ in
\eqref{eq:ns61-average}, then $h((u/n)y)\to h(y)$ strongly,
uniformly for $n-1\leq u\leq n$. Summation on $N<n\leq2N$
and $N^2\sum n^{-3}\to3/8$ prove the trace asymptotic.
\end{proof}

The associated exact Gram identity, which retains all off-diagonal
entries, is
\begin{equation}
 \langle\epsilon_m,\epsilon_n\rangle
 =\frac1{mn}\int_{m-1}^{m}\int_{n-1}^{n}
        \langle\rho_u,\rho_v\rangle\,du\,dv.
 \label{eq:ns61-averaged-full-gram}
\end{equation}
Thus smoothing converts the singular adjacent-difference budget into
an average of the original complete Gram kernel. It does not diagonalize
that kernel or remove the old-space projection.

\begin{proposition}[Complete smoothed gain and its still-missing numerator]
\label{prop:ns61-gain}
Form the old Gram $G^{(2)}$, old/new block $C^{(2)}$, new Gram
$D^{(2)}$, and target loads $b^{(2)}$ using $\sigma_n$ and $\tau$.
Set
\[
 c=(G^{(2)})^{-1}b^{(2)},\quad
 h^{(2)}=b_{\rm new}^{(2)}-(C^{(2)})^*c,\quad
 S^{(2)}=D^{(2)}-(C^{(2)})^*(G^{(2)})^{-1}C^{(2)}.
\]
For the same triangular new-coordinate matrix $T_N$ as in
Proposition~\ref{prop:v158-difference-gain}, put
$g^{(2)}=T_N^*h^{(2)}$ and $K_N^{(2)}=T_N^*S^{(2)}T_N$.
Then
\begin{equation}
 0<K_N^{(2)}\leq\mathcal D_N^{(2)},\qquad
 E_N^{(2)}-E_{2N}^{(2)}
 =(g^{(2)})^*(K_N^{(2)})^{-1}g^{(2)}
 \geq\frac{\norm{g^{(2)}}_2^2}{\mathfrak b_N}
 \geq\frac{8N^2}{3\kappa}\norm{g^{(2)}}_2^2.
 \label{eq:ns61-gain}
\end{equation}
Consequently the cofinal estimate
\begin{equation}
 N^2\norm{g_{2^j}^{(2)}}_2^2
 \geq\frac{3\kappa}{8}\alpha_j E_{2^j}^{(2)},
 \quad N=2^j,\quad 0\leq\alpha_j<1,\quad
 \sum_j\alpha_j=\infty
 \label{eq:ns61-open-correlation}
\end{equation}
for all sufficiently large $j$ would prove RH. This lower estimate
is an open input; no necessity claim is made for it.
\end{proposition}
\begin{proof}
The finite Gram matrices are positive definite by injectivity of $J$.
Orthogonal projection, the complete Schur complement and the invertible
coordinate change give the exact identities as in NS-53.
The raw trace budget bounds the largest eigenvalue of the projected
Gram. The resulting contraction
$E_{2^{j+1}}^{(2)}\leq(1-\alpha_j)E_{2^j}^{(2)}$ proves the
conditional conclusion. In particular, the unsmoothed residual
correlation cannot be inserted in~\eqref{eq:ns61-gain}.
\end{proof}

\begin{proposition}[Explicit loads do not supply a projected lower bound]
\label{prop:ns61-loads}
Use the standard convention
$\zeta(1+z)=z^{-1}+\gamma-\gamma_1z+O(z^2)$ for $\gamma_1$.
Then
\begin{equation}
 b_n^{(2)}=\langle\tau,\sigma_n\rangle
 =\frac{P(\log n)}n,
 \quad P(v)=\tfrac12v^2+(2-\gamma)v+3-2\gamma-\gamma_1.
 \label{eq:ns61-loads}
\end{equation}
In difference coordinates, for $n\geq2$,
\[
 \langle\tau,\epsilon_n\rangle
 =\frac{P(\log n)-P(\log(n-1))}{n}
 \sim\frac{\log n+2-\gamma}{n^2}.
\]
For the actual residual the exact expression is instead
\begin{equation}
 g_n^{(2)}=\frac{P(\log n)-P(\log(n-1))}{n}
       -\sum_{m\leq N}c_m\langle\sigma_m,\epsilon_n\rangle,
 \quad N<n\leq2N.
 \label{eq:ns61-retained-cross}
\end{equation}
The sign and size of the second term are not controlled here.
\end{proposition}
\begin{proof}
Since $\sigma_1(x)=1/x$ for $x>1$, the transform
$-\zeta(s)/s^2-1/(1-s)$ is the Mellin transform of the restriction
to $(0,1)$. Its value and derivative at $s=1$ give
\[
 \int_0^1\sigma_1(x)\,dx=2-\gamma,\qquad
 \int_0^1(-\log x)\sigma_1(x)\,dx=3-2\gamma-\gamma_1.
\]
Change variables $u=nx$ in the load integral and split at $u=1$.
The tail contribution is $(\log n)^2/(2n)$, giving
\eqref{eq:ns61-loads}. Taking the adjacent difference and then
subtracting the old projection gives the remaining identities.
The loads agree, with the sign convention for $\sigma_n$ reversed,
with the $q=2$ formula in~\cite{Ehm2024}.
\end{proof}

\paragraph{The exact point at which the lower-bound attempt stops.}
Every new $\epsilon_n$ is supported in $(0,1/N]$, so the numerator
only sees the small-$x$ part of $R_N$. The elementary formula
\[
 A(z)=z-k\log z+\log(k!),\qquad k=\lfloor z\rfloor,
\]
and Stirling's estimate imply, uniformly for $z\geq1$,
$|A(z)-\tfrac12\log(2\pi z)|\leq3/z$; a complete elementary
remainder is given in Section~\ref{sec:ns61-gram-remainder}.
Thus for $0<x\leq1/N$,
\begin{equation}
 \begin{aligned}
 R_N(x)&=a_N\log(1/x)+\ell_N+\mathcal R_N(x),\\
 a_N&=1-\tfrac12\sum_{m\leq N}c_m,\qquad
 \ell_N=\tfrac12\sum_{m\leq N}c_m\log\frac m{2\pi},\\
 |\mathcal R_N(x)|&\leq3x\sum_{m\leq N}m|c_m|.
 \end{aligned}
 \label{eq:ns61-residual-endpoint}
\end{equation}
Writing $\Delta_n f=f(n)-f(n-1)$, the exact first two moments of
$\epsilon_n$ give
\begin{equation}
 g_n^{(2)}=
 \frac{a_N\{\tfrac12\Delta_n(\log^2)+(2-\gamma)\Delta_n\log\}
              +\ell_N\Delta_n\log}{n}
 +\mathcal E_{N,n},
 \quad N<n\leq2N,
 \label{eq:ns61-endpoint-correlation}
\end{equation}
where
\[
 |\mathcal E_{N,n}|\leq
       \frac{\pi^2\sum_{m\leq N}m|c_m|}{8n^2(n-1)}.
\]
The two leading moments follow from $\int h=1$ and
$\int(-\log x)h(x)\,dx=2-\gamma$, obtained by taking the value
and derivative of $-(1-s)\zeta(s)/s^2$ at $s=1$.
For the error, use nonnegativity and the exact first moment
\eqref{eq:ns61-positive-first-moment}. Neither a lower bound for the displayed
coefficient combination relative to $E_N^{(2)}$, nor a small enough
bound for $\sum m|c_m|$, has been obtained. Replacing these optimized
coefficients by the target loads would discard
\eqref{eq:ns61-retained-cross}.

The improvement is therefore a proved elementary denominator and an
explicit description of the remaining arithmetic numerator. It bypasses
the raw unsmoothed norm obstruction by changing both the functions and
the target. It does not bypass the zero obstruction, certify a lower
correlation, or establish convergence. This is the stop condition for
NS-61 unless an independent estimate for the optimized residual is found.
